\clearpage
% Requires: \usepackage{multirow, subcaption, arydshln}
\begin{table}[h!]
\centering \scriptsize
\setlength{\extrarowheight}{0.5pt}
\setlength{\tabcolsep}{3pt}

\begin{subtable}{0.49\linewidth}
\centering
\begin{tabular}{l|lll|lll}
\textbf{Model} & \textbf{Method} & \textbf{CC} & \textbf{Red.} & \textbf{PC} & \textbf{KT} & \textbf{BertS} \\ \hline
\multirow{2}{*}{LLaMa} & BL &  &  & 0.015 & 0.017 & 0.840 \\
 & CV & Center & Avg & 0.730 & 0.563 & 0.803 \\ \cdashline{1-7}[1pt/3pt]
\multirow{2}{*}{Gemma} & BL &  &  & 0.000 & 0.000 & 0.820 \\
 & CV & Center & PCA & 0.063 & 0.023 & 0.813 \\ \cdashline{1-7}[1pt/3pt]
\multirow{2}{*}{Mistral} & BL &  &  & 0.000 & 0.000 & 0.820 \\
 & CV & Diff & PCA & 0.315 & 0.223 & 0.821 \\
\end{tabular}
\caption{eLife}\label{tab:main-results-elife}
\end{subtable}
\hfill
\begin{subtable}{0.49\linewidth}
\centering
\begin{tabular}{l|lll|lll}
\textbf{Model} & \textbf{Method} & \textbf{CC} & \textbf{Red.} & \textbf{PC} & \textbf{KT} & \textbf{BertS} \\ \hline
\multirow{2}{*}{LLaMa} & BL &  &  & 0.055 & 0.019 & 0.831 \\
 & CV & Center & Avg & 0.725 & 0.553 & 0.805 \\ \cdashline{1-7}[1pt/3pt]
\multirow{2}{*}{Gemma} & BL &  &  & -0.002 & -0.003 & 0.833 \\
 & CV & Center & PCA & 0.106 & 0.065 & 0.821 \\ \cdashline{1-7}[1pt/3pt]
\multirow{2}{*}{Mistral} & BL &  &  & 0.000 & 0.000 & 0.833 \\
 & CV & Diff & PCA & 0.213 & 0.189 & 0.830 \\
\end{tabular}
\caption{PLOS}\label{tab:main-results-plos}
\end{subtable}

\vspace{2mm}

\begin{subtable}{0.49\linewidth}
\centering
\begin{tabular}{l|lll|lll}
\textbf{Model} & \textbf{Method} & \textbf{CC} & \textbf{Red.} & \textbf{PC} & \textbf{KT} & \textbf{BertS} \\ \hline
\multirow{2}{*}{LLaMa} & BL &  &  & 0.000 & -0.003 & 0.848 \\
 & CV & Center & Avg & 0.770 & 0.607 & 0.810 \\ \cdashline{1-7}[1pt/3pt]
\multirow{2}{*}{Gemma} & BL &  &  & 0.000 & 0.000 & 0.874 \\
 & CV & Center & PCA & 0.075 & 0.042 & 0.895 \\ \cdashline{1-7}[1pt/3pt]
\multirow{2}{*}{Mistral} & BL &  &  & 0.000 & 0.000 & 0.874 \\
 & CV & Center & PCA & 0.121 & 0.098 & 0.899 \\
\end{tabular}
\caption{Eureka}\label{tab:main-results-eureka}
\end{subtable}
\hfill
\begin{subtable}{0.49\linewidth}
\centering
\begin{tabular}{l|lll|lll}
\textbf{Model} & \textbf{Method} & \textbf{CC} & \textbf{Red.} & \textbf{PC} & \textbf{KT} & \textbf{BertS} \\ \hline
\multirow{2}{*}{LLaMa} & BL &  &  & 0.030 & -0.004 & 0.840 \\
 & CV & Center & Avg & 0.674 & 0.573 & 0.801 \\ \cdashline{1-7}[1pt/3pt]
\multirow{2}{*}{Gemma} & BL &  &  & 0.002 & 0.001 & 0.820 \\
 & CV & Center & Avg & 0.129 & 0.110 & 0.821 \\ \cdashline{1-7}[1pt/3pt]
\multirow{2}{*}{Mistral} & BL &  &  & 0.002 & 0.001 & 0.820 \\
 & CV & Center & Avg & 0.141 & 0.177 & 0.809 \\
\end{tabular}
\caption{SciNews}\label{tab:main-results-scinews}
\end{subtable}

\caption[Correlation of the specified level of readability with the Coleman-Liau readability index.]{Pearson Correlation (PC) and Kendall-Tau (KT) correlation of the specified level of readability with the Coleman-Liau readability index, and BertScore F1 (BertS). Reported is the best baseline (BL). We report the best performing CV setting for Contrastive Combination (CC) and Reduction (Red.)}\label{tab:ch-4-main-results-mini}
\end{table}
\section{Results}\label{sec:4-chapter-results}
We report best settings per dataset in \Cref{tab:ch-4-main-results-mini} and full ablations in~\Cref{appendix:ch-4-full-results}. 

We see that Avg works the best as a reduction setting for all datasets except SciNews. This holds for both Diff and Center as the Contrastive Combination setting. For eLife, PLOS, and Eureka, Center achieves slightly better performance, although the difference is small. In general, the choice of reduction method makes the largest difference in terms of performance. While the Center-Avg settings perform best for 3 of the 4 datasets, SciNews achieving better performance with Diff-PCA indicates this setting is not universally optimal. Some tuning of the control vector settings may be necessary, although the overhead of this tuning is minimal compared to that of higher compute methods, such as finetuning, prompt optimization, and reinforcement-learning.

The failure of the training methods does not mean they did not learn how to summarize but rather they were unable to learn to control the readability of the summaries.  ICL, SFT, QLoRA/LoRA SFT, and prompt optimization perform poorly, with ICL and prompt optimization aligning with prior work showing they are insufficient for readability control \cite{ribeiro-etal-2023-generating}. Training-based methods also underperform, likely due to dataset differences (scientific texts vs. shorter news and fewer examples). Example outputs appear in Tables \ref{tab:sft-examples}–\ref{tab:cv-ex-outputs}. We include a more detailed discussion of the baseline performance in~\Cref{appendix:chap-4-baseline-discussion}.


For most settings, using Center for the Contrastive Combination and Avg for the Reduction step performs best. Additionally, although the control vectors use eLife for the extraction dataset, the best control vectors generalize well to the other datasets, achieving similarly high correlations. This suggests that the resulting control vectors are in fact representing readability, rather than dataset specific information. 


\subsection{Analysis of extraction dataset requirements.}\label{sec:other-datasets-exp}
We conduct experiments using the non-compliant datasets. We extract control vectors using the 128 most readable summaries from each train dataset, similar to the process in~\S\ref{sec:data-selection}. For evaluation, we randomly sample 248 examples from each test set. At inference, we use the control vectors extracted from each dataset's respective extraction set. This is an easier task than that presented above as it does not test generalization to other datasets. We use Center as our Contrastive Combination method and Avg as our reduction method, the best performing setting above. Results are below:
\input{research-chapters/4-chapter/tables/other-dataset-results}

We find that the non-compliant datasets perform poorly. Eureka results in a moderate negative correlation, likely a result of the extraneous information present in the lay summaries, which can increase the readability scores. SciNews and PLOS have low correlation scores, indicating poor controllability. This, contrasted with the high performance presented in~\Cref{tab:ch-4-main-results-mini}, provides evidence for the validity of our requirement framework.

\begin{table*}[ht]
\centering
\footnotesize
\begin{tabular}{ll|llll|llll}
\multicolumn{2}{c|}{\textbf{Settings}}&
\multicolumn{4}{c|}{\textbf{Per Sample Latency (ms)}} &
\multicolumn{4}{c}{\textbf{Per Token Latency (ms)}} \\
BSZ & Seq Len &
Std & CV & $\Delta$ & \% OH &
Std & CV & $\Delta$ & \% OH \\
\hline
1  & 512  & 46.123 & 48.434 & 2.311 & 5.01 &
0.090 & 0.095 & 0.005 & 5.01 \\
1  & 1024 & 85.334 & 89.729 & 4.395 & 5.15 &
0.083 & 0.088 & 0.004 & 5.15 \\
1  & 2048 & 163.423 & 172.033 & 8.610 & 5.27 &
0.080 & 0.084 & 0.004 & 5.27 \\
\hline
4  & 512  & 39.693 & 41.864 & 2.171 & 5.47 &
0.078 & 0.082 & 0.004 & 5.47 \\
4  & 1024 & 78.111 & 82.373 & 4.262 & 5.46 &
0.076 & 0.080 & 0.004 & 5.46 \\
4  & 2048 & 156.692 & 165.435 & 8.743 & 5.58 &
0.077 & 0.081 & 0.004 & 5.58 \\
\hline
16 & 512  & 38.098 & 40.223 & 2.125 & 5.58 &
0.074 & 0.079 & 0.004 & 5.58 \\
16 & 1024 & 75.600 & 79.912 & 4.312 & 5.70 &
0.074 & 0.078 & 0.004 & 5.70 \\
16 & 2048 & 155.383 & 163.905 & 8.522 & 5.48 &
0.076 & 0.080 & 0.004 & 5.48 \\
\end{tabular}
\caption[Latency comparison between inference on a Standard LLM and an LLM with added Control Vectors.]{Latency comparison between inference on a Standard (Std) LLM and an LLM with added Control Vectors (CV). We report the per sample and per token latency in milliseconds (ms), the $\Delta$ latency from using control vectors, and the \% overhead (OH).}
\label{tbl:latency-results}
\end{table*}

\subsection{Latency testing}\label{appendix:latency-testing}

We conduct latency tests on the added overhead of using control vectors over standard LM inference. We randomly generate input IDs and measure the latency of the forward pass, as this is the portion of the generation pipeline where control vectors are used and have the potential to add overhead. We experiment with 3 batch sizes (1,4, 16) and 3 sequence lengths (512,1024,2048) for a total of 9 ablations. For each setting, we run 32 warm-up runs then measure the latency for 1024 inference calls. We report the mean latency in milliseconds per sample and per token. We additionally report the $\Delta$ latency and percent overhead. All the latency tests are run using Llama 3.1 8b Instruct (\texttt{meta-llama/Llama-3.1-8B-Instruct}) on a single NVIDIA A100 GPU, using half-precision floating-point. The results are reported in~\Cref{tbl:latency-results}. We find that using control vectors adds approximately a 5\% latency overhead when measured per sample or per token. We believe this difference is negligible when compared to other controllable generation methods, that require expensive training or hyper-parameter searches, or the additional token processing required for In-Context Learning. 